\bookchapter{第六十一章：Swift 速成——蘋果開發之第一語言}

\chapterintro{此章吾儕正式學其語言。勿懼——Swift 素以「善待新手」著稱，且爾不必背語法之書：先掌握 20\% 之語法，即可看懂且寫出 80\% 日常之代碼。其餘者，且作項目且查之。記吾儕學習之則：看懂勝於強記，會用勝於精通。}

\booksection{一、變量與類型：為數據起名}

程序者，「數據加邏輯」也。變量者，為數據起一名而儲之也。Swift 以 let 與 var 聲明之：

\begin{lstlisting}[language=Swift]
let name = "張三"        // let：常量（宜以為默認）
var age = 18             // var：變量（後尚可改）

age = 19                 // <i class="fas fa-check-circle mr-1"></i> var 可改
// name = "李四"         // <i class="fas fa-times-circle mr-1"></i> let 不可改，編譯報錯

let score: Int = 100     // 亦可明書其型
let pi = 3.14            // 不書亦自推之

print("汝好，\(name)")   // 字符串插值：反斜槓 + 括號
\end{lstlisting}

常用類型一表記之：

\begin{center}
\begin{tabularx}{\linewidth}{|X|X|X|}
\hline
\textbf{類型} & \textbf{儲何物} & \textbf{示例} \tabularnewline
\hline
Int & 整數 & 42 \tabularnewline
Double & 小數 & 3.14 \tabularnewline
String & 文字 & 「你好」 \tabularnewline
Bool & 真/假 & true / false \tabularnewline
{}[Int]（陣列） & 一組數據 & {}[1, 2, 3] \tabularnewline
\hline
\end{tabularx}
\end{center}

\begin{center}\small\color{gray}Swift 常用基礎類型\end{center}

\begin{tcbinfo}{小訣：let 優先之原則}
Swift 社區之鐵律：默認用 let，確有所須方用 var。可用 let 之處而用 var，Xcode 將以黃字警爾。此習能助爾少卻一大類「變量遭意外之改」之 bug——代碼之變愈少，推理愈易。

\end{tcbinfo}

\booksection{二、可選值（Optional）：Swift「防崩護身之符」}

此 Swift 至要之概念，無出其右。一言以蔽之：一值或「不存」之時，Swift 強爾理之——自語言之層滅「空指針崩潰」此業中之惡夢。

\begin{lstlisting}[language=Swift]
var name: String = "張三"   // 普通之型：決不容為 nil
// name = nil                // <i class="fas fa-times-circle mr-1"></i> 編譯之誤，徑阻之

var nickname: String? = nil  // 加問號：許為 nil（可選值）
nickname = "三兒"

// 用可選值，必"解包"：
if let n = nickname {        // if let：安全解包
    print("汝好，\(n)")
} else {
    print("尚無暱稱")
}

let len = nickname?.count ?? 0   // ?? 予默認值（空合運算符）
print(len)
\end{lstlisting}

\begin{tcbwarn}{注意：! 乃「高危操作」}
代碼中書 nickname!（強制解包）猶謂編譯器曰「吾保其非 nil」——一猜而誤，App 當場即崩。新手之期其則曰：務用 if let / ??，勿用 !。Xcode 補全代碼之時偶爾順手加 !，記而易之。

\end{tcbwarn}

\booksection{三、函數與閉包：打包邏輯成「鈕」}

函數者，一段有名之邏輯也，若遙控之鈕。閉包（Closure）者，「無名之函」也，可隨手遞於人：

\begin{lstlisting}[language=Swift]
// 函數：名字於前，參數與返回類型於後
func greet(name: String) -> String {
    return "汝好，\(name)！"
}
print(greet(name: "張三"))     // 調用時要帶參數名

// 閉包：就地定義，輕快直接
let numbers = [1, 2, 3, 4]
let doubled = numbers.map { $0 * 2 }   // $0 = 當前元素
// [2, 4, 6, 8]

let big = numbers.filter { $0 > 2 }    // 篩選
// [3, 4]
\end{lstlisting}

SwiftUI 之中閉包無所不在：按鈕 Button(「點我」) \{ ... \} 之後大括號即一閉包也。map、filter、sorted 此數「集合連招」請牢記之——來日 90\% 列表之處理皆其組合也。

\booksection{四、結構體與協議：Swift「積木」之體系}

Swift 官方之好甚明：先用結構體（struct），須繼承之時方用類（class）。結構體者「值類型」也——複製之時各自為之，天然少一類之 bug：

\begin{lstlisting}[language=Swift]
// 結構體：盛數據之盒
struct Record {
    var id = UUID()
    var amount: Double
    var category: String
    var date: Date

    // 結構體自帶"逐一成員構造器"：可徑建
}
let r = Record(amount: 25.5, category: "餐飲", date: Date())

// 協議：約"必具何能"
protocol Identifiable {
    var id: UUID { get }
}
struct Record: Identifiable { ... }   // 遵守協議 = 承諾實現其要求
\end{lstlisting}

SwiftUI 之中幾乎皆 struct（View 即協議也）。此言今不解無妨，第 62 章將日日相見——Swift 之世界，先擁抱 struct 可也。

\booksection{五、貽零基者之語法速查卡}

\begin{center}
\begin{tabularx}{\linewidth}{|X|X|}
\hline
\textbf{欲何為} & \textbf{寫法} \tabularnewline
\hline
定常量/變量 & \texttt{let x = 1}/\texttt{var x = 1} \tabularnewline
字符串插值 & 「你好，\textbackslash{}(name)」 \tabularnewline
定函數 & func 名(參數: 類型) -> 返回類型 \{ ... \} \tabularnewline
判斷 & if 條件 \{ ... \} else \{ ... \} \tabularnewline
遍歷 & for x in 陣列 \{ ... \} \tabularnewline
安全解包 & if let x = 可選值 \{ ... \} \tabularnewline
默認值 & x ?? 默認值 \tabularnewline
定結構體 & \texttt{struct X \{ var a: Int; var b: String \}} \tabularnewline
閉包 & \texttt{\{ \$0 * 2 \}} \tabularnewline
\hline
\end{tabularx}
\end{center}

\begin{center}\small\color{gray}Swift 語法速查卡（新手版）\end{center}

\begin{tcbpractice}{來動手：於 Playgrounds 玩 15 分鐘之 Swift}
有 Mac 則開 Xcode Playground；無 Mac 則用 iPad 上之 Swift Playgrounds 或網頁版之演練場。以本章每一段之代碼皆跑一遍，改其數字與文字，故意造數誤以觀 Xcode 何以提示之。此「先玩後學」之 15 分鐘，將令爾第 62 章上手之時輕鬆一倍。

\end{tcbpractice}

\begin{tcbsummary}{本章小結}
\begin{itemize}
  \item let 為先，var 慎用；類型自動推斷，冒號之後明言亦可
  \item 可選值者 Swift 之魂也：或空必理之，if let 與 ?? 乃護身之符，! 乃高危操作
  \item 函數之調用須帶參數之名；閉包者無名之函也，SwiftUI 之中無所不在
  \item 先用 struct，守協議（Protocol）；SwiftUI 之中幾乎皆 struct
  \item 20\% 之語法為 80\% 之事：先會用之，再精通之
\end{itemize}

\end{tcbsummary}
